\section{Introduction}
\label{sec:intro}

Bug reports are critical input to the software development process. However, without the ability to reproduce a bug, developers cannot diagnose or fix it. Recent work shows that fewer than 8\% of bug reports on platforms like Jira meet the criteria for first-time reproducibility, creating a significant bottleneck in software maintenance. For large projects managing tens of thousands of reports monthly (such as TensorFlow or Django), automating the reproducibility evaluation step offers clear practical value.

Large language models (LLMs) like GPT-4o have shown promise in various software engineering tasks, including bug report analysis and improvement. Yet a critical gap remains: no prior work has measured whether an LLM can evaluate bug report reproducibility with sufficient agreement to developer judgment. Existing studies either use LLMs to generate or improve bug reports, or measure agreement between human raters alone---not between LLM evaluators and developer assessment.

This work addresses that gap by investigating whether an LLM-based reproducibility evaluator can achieve substantial agreement (Cohen's $\kappa \geq 0.70$) with manual human review. We use the ImproBR replication package~\cite{akyol2026improbr}, which is uniquely suited to the question because it contains two forms of the same 139 Mojira Minecraft bug reports: the original \emph{Raw} reports as submitted, and \emph{Improved} versions in which an LLM pipeline has rewritten each report to be more complete. Both forms were independently annotated by two human reviewers using the same rubric.

We treat the two forms as \emph{co-equal conditions} and run the identical evaluation on each. This lets us ask not only whether GPT-4o mini agrees with human reproducibility judgment on ordinary reports, but also whether that agreement holds up on reports that another LLM has already polished. The two conditions are complementary: the Raw condition tests the model on the messy text a triage system actually receives, while the Improved condition tests whether the model can still track human judgment when the surface text has been made fluent and complete-looking by AI. A large gap between the two would be a warning that an LLM evaluator can be misled by fluent presentation rather than genuine reproducibility.

The contribution is three-fold: we establish a reusable measurement protocol for LLM-as-evaluator agreement on bug report quality; we provide an empirical baseline for whether LLMs can reliably assist in bug triage on both original and AI-improved reports; and we show that the two answers differ, which has direct implications for anyone considering an LLM evaluator in a pipeline that already contains LLM-generated text.
